\documentclass[letterpaper,12pt, twocolumn]{report}

\usepackage[a4paper, margin=0.8in]{geometry} 

\setlength{\columnsep}{0.5cm} % increase column gap
\begin{document}

% title page
\begin{titlepage}
    \begin{center}
        \LARGE
        \textbf{Stock Movement Prediction from Price Trends and News Sentiment}

        % \LARGE
        % \vspace{5pt}
        % ECS 171

        \normalsize

        \vspace{40pt}
        \textbf{Team}: Honoré Alexander, Owen Holt, Pranavi Khanna,  Dylan Lim, Yihong Li, Ethan Lee, Dan Firstenberg, Hyeongseung Nam, Oscar Pineda, Kevin Zhang, Zachary Chan, Vicente Aguayo\\

        \vspace{25pt}
        November 20th, 2025
    
        \vspace{40pt}
        \textbf{Abstract} \\
        \vspace{10pt}

    \end{center}
        \raggedright
   Within recent years, machine learning has emerged as one of the most powerful tools in predicting stock performances for investors. Such machine learning algorithms integrate market data, financial reports, and macroeconomic data to predict stock movement and volatility; however leading methods systematically overreact to macroeconomic news or fail to take it into account completely. A central aim of our project is to quantify how exchange-traded fund (ETF) relationships reflect underlying macroeconomic conditions, and how these relationships influence individual stock behavior. To address this limitation, we constructed a data management process linking historical company-specific news headlines to corresponding stock performance data across 60 individual tickers.
   
   (Add 1-2 more sentences for method once done with ML side) 
(Add 2-3 more sentences on results)
  

\vfill

% Horizontal line and keywords at bottom
\hrule
\vspace{0.5cm}
\noindent
\textbf{Keywords:} Machine Learning, Exchange-Traded Funds(ETF’s), Batch Gradient Descent, Ordinary Least Squares 



\end{titlepage}

% report page 1 begins
\newpage


\large
\noindent
\textbf{1 Introduction}\\
\normalsize



\large
\noindent
\textbf{2 Methods} \\
\normalsize
To train our model, we followed a conventional machine learning development model.

\noindent
\textbf{2.1 Data Sourcing and Preprocessing}\\
\textit{2.1.1 Dataset Identification}\\
text here \\
\textit{2.1.1 Dataset Merging} \\


\large
\noindent
\textbf{3 Results} \\
\normalsize

\large
\noindent
\textbf{4 Discussion} \\
\normalsize

\large 
\noindent
\textbf{5 Conclusion} \\
\normalsize


\newpage
\large
\noindent
\textbf{References} 
\normalsize
\raggedright
\begin{enumerate}
    \item Usmani, Shazia, and Jawwad A Shamsi. “LSTM based stock prediction using weighted and categorized financial news.” PloS one vol. 18,3 e0282234. 7 Mar. 2023, 
    \item Weijia (Vivian) Li \& Kara M. Kockelman, 2022. "How does machine learning compare to conventional econometrics for transport data sets? A test of ML versus MLE," Growth and Change, Wiley Blackwell, vol. 53(1), pages 342-376, March.
    \item Saberironaghi, M.; Ren, J.; Saberironaghi, A. Stock Market Prediction Using Machine Learning and Deep Learning Techniques: A Review. AppliedMath 2025, 5, 76. 
    \item Suresh, Harini, and John Guttag. “A Framework for Understanding Sources of Harm Throughout.” Arxiv.Org, 2021,
    \item Chen, Qi, et al. “Price Informativeness and Investment Sensitivity to Stock Price Qi Chen.” Finance.Wharton.Upenn.Edu, July 2006, 
    \item Malanin, V. Comparative study of neural network architectures in medical survival analysis: Evaluation of MLP, CNN, and LSTM models. International Journal of Neural Networks and Deep Learning, 2(2), 24–43, 2025, 
\end{enumerate} 


\newpage
\vspace{20pt}
\large
\noindent
\textbf{Author Contributions} 
\normalsize
\begin{itemize}
    \item \textbf{Honoré Alexander}: Led backend development of neural network system, handling embeddings, deriving price features, and optimizing model to predict if the stock will close higher or lower.
    \item \textbf{Owen Holt}
    \item \textbf{Pranavi Khanna}
    \item \textbf{Dylan Lim}: Created application and backend for model inference, designed the initial frontend. Other contributions to source research and technical report documentation.
    \item \textbf{Yihong Li}: Researched sources and drafted most of the technical report sections. Set up LaTeX and formatted the report. Set up and designed frontend for model inference.
    \item \textbf{Ethan Lee}
    \item \textbf{Dan Firstenberg}
    \item \textbf{Hyeongseung Nam}
    \item \textbf{Oscar Pineda}: Researched sources for the report, contributed to writing the report
    \item \textbf{Kevin Zhang}
    \item \textbf{Zachary Chan}
    \item \textbf{Vicente Aguayo}
\end{itemize}




\end{document}


